\section*{Reproducibility}
\addcontentsline{toc}{section}{Reproducibility}

Every number in this paper is produced by a committed script, and no figure or
table was edited by hand. The four scripts write the files the paper
\verb|\input|s, so a stale table cannot survive a regeneration unnoticed.

\begin{center}
\small
\begin{tabular}{lll}
\toprule
Script & Produces & Section \\
\midrule
\verb|experiment_qp_identities.py| & \verb|quadprog_identities*.tex|
  & \S\ref{ssec:identities} \\
\verb|experiment_qp_pdas.py| & \verb|quadprog_pdas.pdf|, \verb|_defs.tex|
  & \S\ref{ssec:pdas} \\
\verb|experiment_qp_sweep.py| & \verb|quadprog_sweep.pdf|, \verb|_defs.tex|
  & \S\ref{ssec:sweepcost} \\
\verb|experiment_qp_compare.py| & \verb|quadprog_compare*.tex|, \verb|.pdf|
  & \S\ref{ssec:compare} \\
\bottomrule
\end{tabular}
\end{center}

\paragraph{Software under study.} \texttt{cvx-quadprog}, version
$0.4.1$~\cite{cvxquadprog}, archived with a DOI. The measurements are tied to that
version rather than to the latest release, so that the reference keeps pointing at
the code that produced them. The experiments import the package's private modules
where the claims are about private state --- $J$, the packed $R$, and the vectors
$d$, $z$, $r$ --- and instrument the shipped routines by wrapping them rather than
reimplementing them, so what is measured is what runs.

\paragraph{Regenerating everything.} From the paper's source tree, \verb|make figures|
runs all four scripts and rewrites every generated file; \verb|make compile| then
rebuilds this document. Each script also runs standalone as a module. Setting
\verb|EXPERIMENT_SMOKE=1| shrinks every sweep so that the whole set completes in
seconds, which is how the continuous-integration suite exercises each code path
without reproducing the full runs, and \verb|EXPERIMENT_OUT| redirects the outputs
so that a reduced run cannot overwrite the committed ones.

\paragraph{Determinism.} Every random problem is drawn from a seed derived
arithmetically from the family, size and instance index, so the tables reproduce
across runs and across machines. Timings do not, and are not claimed to. The
figures in \S\ref{ssec:sweepcost} and \S\ref{ssec:compare} were taken on one
machine, against NumPy built on Apple Accelerate, with no thread count set. That
last detail is not incidental: this solver pushes its work into BLAS calls, so its
timings inherit whatever threading decision the installed library makes by
default, and those defaults differ between libraries by more than the effects
being measured here. Accelerate exposes no thread control, so there was nothing to
fix; on a library that does, the same experiment can read very differently. The
timings should be re-measured rather than quoted. The residuals in \S\ref{ssec:identities} and the counts in
\S\ref{ssec:pdas} are machine-independent up to the last digit or two of
floating-point rounding, and are the numbers this paper actually rests on.

\paragraph{Optional dependency.} The comparison in \S\ref{ssec:compare} includes
the reference C implementation of the same algorithm where a wheel for it is
installed. Building it needs the C toolchain that the package under study exists
to avoid, so the script treats the import as optional and omits that row, with a
printed note, on a host that lacks it. Every other row is unconditional.
